\documentclass[12pt, a4paper]{article}

% ─── Packages ────────────────────────────────────────────────────────────────
\usepackage[margin=2.5cm]{geometry}
\usepackage{amsmath, amssymb, amsthm}
\usepackage{graphicx}
\usepackage{booktabs}
\usepackage{array}
\usepackage{longtable}
\usepackage{xcolor}
\usepackage{hyperref}
\usepackage{natbib}
\usepackage{caption}
\usepackage{subcaption}
\usepackage{listings}
\usepackage{float}
\usepackage{setspace}
\usepackage{titlesec}
\usepackage{fancyhdr}
\usepackage{enumitem}
\usepackage{tcolorbox}
\usepackage{mdframed}
\usepackage{pgfplots}
\pgfplotsset{compat=1.18}

% ─── Colours ─────────────────────────────────────────────────────────────────
\definecolor{jpmblue}{RGB}{0, 44, 95}
\definecolor{jpmgold}{RGB}{180, 150, 80}
\definecolor{lightgray}{RGB}{245, 245, 245}
\definecolor{codegray}{RGB}{240, 240, 240}
\definecolor{placeholder}{RGB}{200, 50, 50}
\definecolor{linkblue}{RGB}{0, 90, 170}

% ─── Hyperref ────────────────────────────────────────────────────────────────
\hypersetup{
    colorlinks=true,
    linkcolor=linkblue,
    citecolor=jpmblue,
    urlcolor=linkblue,
    pdftitle={Term Loan Pricing Model -- JP Morgan Interview Project},
    pdfauthor={Candidate},
}

% ─── Page style ──────────────────────────────────────────────────────────────
\pagestyle{fancy}
\fancyhf{}
\fancyhead[L]{\small\color{jpmblue}\textbf{JP Morgan -- Term Loan Pricing Model}}
\fancyhead[R]{\small\color{jpmblue}\today}
\fancyfoot[C]{\thepage}
\renewcommand{\headrulewidth}{0.4pt}
\renewcommand{\headrule}{\hbox to\headwidth{\color{jpmblue}\leaders\hrule height \headrulewidth\hfill}}

% ─── Section styling ─────────────────────────────────────────────────────────
\titleformat{\section}{\large\bfseries\color{jpmblue}}{}{0em}{\thesection\quad}[\titlerule]
\titleformat{\subsection}{\normalsize\bfseries\color{jpmblue}}{}{0em}{\thesubsection\quad}
\titleformat{\subsubsection}{\normalsize\itshape\color{jpmblue}}{}{0em}{\thesubsubsection\quad}

% ─── Code listing style ──────────────────────────────────────────────────────
\lstdefinestyle{python}{
    language=Python,
    backgroundcolor=\color{codegray},
    basicstyle=\ttfamily\small,
    keywordstyle=\color{jpmblue}\bfseries,
    stringstyle=\color{jpmgold},
    commentstyle=\color{gray}\itshape,
    numberstyle=\tiny\color{gray},
    numbers=left,
    stepnumber=1,
    frame=single,
    rulecolor=\color{jpmblue!30},
    breaklines=true,
    captionpos=b,
    showstringspaces=false,
    tabsize=4,
}
\lstset{style=python}

% ─── Placeholder environment ─────────────────────────────────────────────────
\newmdenv[
  backgroundcolor=red!8,
  linecolor=placeholder,
  linewidth=1.5pt,
  roundcorner=4pt,
  innerleftmargin=10pt,
  innerrightmargin=10pt,
  innertopmargin=8pt,
  innerbottommargin=8pt,
]{placeholder}

\newcommand{\phtext}[1]{%
  \begin{placeholder}
    \color{placeholder}\textbf{[PLACEHOLDER]}~#1
  \end{placeholder}
}

% ─── Theorem environments ────────────────────────────────────────────────────
\theoremstyle{definition}
\newtheorem{definition}{Definition}[section]
\newtheorem{remark}{Remark}[section]

% ─── Spacing ─────────────────────────────────────────────────────────────────
\setstretch{1.15}
\setlength{\parskip}{6pt}

% ─── Title ───────────────────────────────────────────────────────────────────
\title{
  \vspace{-1cm}
  {\color{jpmblue}\rule{\linewidth}{2pt}}\\[0.4cm]
  {\Large\bfseries\color{jpmblue} Term Loan Interest Rate Pricing Model}\\[0.2cm]
  {\large\color{jpmgold} JP Morgan Interview Project}\\[0.3cm]
  {\color{jpmblue}\rule{\linewidth}{0.8pt}}
}
\author{
  \textbf{Candidate} \\
  \small Quantitative Finance $\cdot$ Credit Risk Modelling
}
\date{\today}

% ═════════════════════════════════════════════════════════════════════════════
\begin{document}
% ═════════════════════════════════════════════════════════════════════════════

\maketitle
\thispagestyle{fancy}

\begin{abstract}
\noindent
This report presents a quantitative framework for pricing the interest rate of term loans.
Term loan pricing is decomposed into a risk-free component, proxied by the US Treasury yield
of the same maturity, and a credit spread that compensates the lender for default risk,
recovery uncertainty, and illiquidity.
We survey the academic literature on defaultable bond and loan pricing, propose a
two-stage machine-learning model, discuss its applicability to private and unlisted
borrowers, extend the framework to one-month resale price forecasting, and demonstrate
how to compute 95\% confidence intervals around both the static price and the
one-month forecast.
Code and data used in this study are available at:
\begin{center}
  \href{https://github.com/PLACEHOLDER/loan-pricing}{\texttt{https://github.com/PLACEHOLDER/loan-pricing}}
\end{center}
\end{abstract}

\tableofcontents
\newpage

% ─────────────────────────────────────────────────────────────────────────────
\section{Introduction}
% ─────────────────────────────────────────────────────────────────────────────

A \textbf{term loan} is a debt instrument in which the borrower pays a fixed periodic
coupon and repays the principal in full at maturity.
The all-in interest rate charged by the bank consists of two additive components:

\begin{equation}
  r_{\text{loan}} = r_f(T) + s,
  \label{eq:allin}
\end{equation}

\noindent where $r_f(T)$ is the yield of the US Treasury bond with the same maturity $T$,
and $s > 0$ is the \emph{credit spread}, which compensates the lender for
(i) the probability that the borrower defaults (\textit{PD}),
(ii) the loss given default (\textit{LGD}), and
(iii) the illiquidity of the loan instrument.

The objective of this project is to build a data-driven model that recommends the
appropriate credit spread $s$ for a new loan, based on borrower characteristics,
loan terms, and prevailing market conditions.

% ─────────────────────────────────────────────────────────────────────────────
\section{Model Choice and Literature Survey}
\label{sec:literature}
% ─────────────────────────────────────────────────────────────────────────────

\subsection{Overview of the Pricing Problem}

Under the risk-neutral measure, the credit spread of a defaultable zero-coupon bond
(and by extension a term loan) can be expressed as
\citep{duffie1999modeling}:

\begin{equation}
  s \approx \lambda \cdot \text{LGD},
  \label{eq:duffiesingelton}
\end{equation}

\noindent where $\lambda$ is the risk-neutral default intensity and $\text{LGD} = 1 - R$,
with $R$ the recovery rate.
This fundamental relationship motivates the two-stage modelling strategy adopted in
Section~\ref{sec:architecture}: first estimate $\lambda$ (equivalently, the
probability of default), then map $(\lambda, \text{LGD})$ to a market spread.

\subsection{Structural Models}

\textbf{Structural models} treat default as an endogenous economic event triggered
when the market value of the firm's assets falls below a threshold.

\begin{table}[H]
\centering
\caption{Key structural models for credit spread pricing.}
\label{tab:structural}
\begin{tabular}{llp{7cm}}
\toprule
\textbf{Model} & \textbf{Year} & \textbf{Key Contribution} \\
\midrule
\citet{merton1974pricing}   & 1974 & Equity is a call option on firm assets; default at maturity when $V_T < D$. \\
\citet{blackcox1976}        & 1976 & First-passage time; default can occur before maturity. \\
\citet{longstaffschwartz1995} & 1995 & Stochastic interest rates with mean reversion. \\
\citet{leland1994}          & 1994 & Endogenous leverage and default trigger; tax shields. \\
KMV / Moody's Analytics     & 1990s & Distance-to-Default $\to$ Expected Default Frequency (EDF). \\
\bottomrule
\end{tabular}
\end{table}

The Merton (1974) model defines the distance to default as
\begin{equation}
  \text{DD} = \frac{\ln(V_A/D) + (\mu - \tfrac{1}{2}\sigma_A^2)T}{\sigma_A \sqrt{T}},
\end{equation}
with $V_A$ the asset value, $D$ the face value of debt, $\mu$ the drift, and
$\sigma_A$ the asset volatility.
\textbf{Limitation:} these models require observable equity prices and implied
volatilities, making them inapplicable to private borrowers.

\subsection{Reduced-Form (Intensity) Models}

\textbf{Reduced-form models} treat the default event as an exogenous Poisson process
governed by a hazard rate (intensity) $\lambda_t$.
The price of a defaultable zero-coupon bond under the risk-neutral measure is
\begin{equation}
  P(t, T) = \mathbb{E}^Q\!\left[\exp\!\left(-\int_t^T (r_s + \lambda_s \cdot \text{LGD})\,ds\right)\right].
\end{equation}

\begin{table}[H]
\centering
\caption{Key reduced-form models.}
\label{tab:reducedform}
\begin{tabular}{llp{7.5cm}}
\toprule
\textbf{Model} & \textbf{Year} & \textbf{Key Contribution} \\
\midrule
\citet{jarrowturnbull1995}  & 1995 & Default = Poisson process; constant intensity. \\
\citet{duffie1999modeling}  & 1999 & Stochastic $\lambda$; spread $\approx \lambda \cdot \text{LGD}$ in closed form. \\
\citet{jarrowlando1997}     & 1997 & Rating-based Markov chain transition model. \\
\citet{lando1998cox}        & 1998 & Cox (doubly stochastic) process; $\lambda$ driven by state variables. \\
\bottomrule
\end{tabular}
\end{table}

\subsection{Machine Learning and Empirical Models}

Recent empirical work has demonstrated that accounting-based and hybrid
ML models achieve competitive or superior predictive accuracy compared to
structural and reduced-form models, especially for cross-sectional pricing
\citep{altman1968financial, mai2019deep}.

\begin{itemize}[leftmargin=*]
  \item \textbf{Altman Z-Score \citep{altman1968financial}:}
    A linear discriminant function of five financial ratios, serving as the
    canonical accounting-based default predictor.
  \item \textbf{Logistic Regression:}
    Regress default indicator on financial ratios; output is a calibrated PD.
  \item \textbf{Gradient Boosting (XGBoost / LightGBM):}
    State-of-the-art for tabular financial data; handles non-linearities and
    interactions automatically \citep{chen2016xgboost}.
  \item \textbf{Deep Learning \citep{mai2019deep}:}
    Multi-layer perceptrons and LSTMs applied to sequential financial statements.
  \item \textbf{Conformal Prediction:}
    Model-agnostic framework for constructing coverage-guaranteed prediction
    intervals \citep{angelopoulos2023conformal}.
\end{itemize}

\subsection{Chosen Model Architecture}
\label{sec:architecture}

We adopt a \textbf{two-stage gradient-boosted tree model}:

\begin{enumerate}
  \item \textbf{Stage 1 — PD Estimation:} An XGBoost classifier trained on borrower
    financial ratios outputs the probability of default $\widehat{\text{PD}}$.
  \item \textbf{Stage 2 — Spread Regression:} An XGBoost regressor maps
    $(\widehat{\text{PD}},\,\text{loan features},\,\text{macro variables})$ to the
    credit spread $\hat{s}$.
\end{enumerate}

The all-in rate is then computed via Equation~\eqref{eq:allin}.
This architecture works for both public and private borrowers (see Section~\ref{sec:private})
and readily admits uncertainty quantification via quantile regression
(see Section~\ref{sec:ci}).

% ─────────────────────────────────────────────────────────────────────────────
\section{Data Sources and Features}
\label{sec:data}
% ─────────────────────────────────────────────────────────────────────────────

\subsection{Data Sources}

\begin{table}[H]
\centering
\caption{Data sources used in model training.}
\label{tab:data}
\begin{tabular}{p{4cm} p{5.5cm} p{3.5cm}}
\toprule
\textbf{Source} & \textbf{Content} & \textbf{Access} \\
\midrule
SEC EDGAR XBRL API       & Standardised financial statements (10-K/10-Q) & Free \\
FRED (St.\ Louis Fed)    & Treasury yields, macro indices, IG/HY OAS      & Free \\
WRDS DealScan            & Historical syndicated loan spreads              & Academic \\
SBA 7(a) Loan Data       & Loan outcomes for SME borrowers                 & Free \\
ICE BofA Indices (FRED)  & Investment-grade / high-yield OAS               & Free \\
Compustat                & Standardised financial ratios                   & WRDS/Paid \\
\bottomrule
\end{tabular}
\end{table}

\subsection{Feature Set}

\begin{table}[H]
\centering
\caption{Model features grouped by category.}
\label{tab:features}
\begin{tabular}{p{3.5cm} p{4cm} p{4cm}}
\toprule
\textbf{Borrower Financials} & \textbf{Loan Characteristics} & \textbf{Market / Macro} \\
\midrule
Debt / EBITDA & Maturity (years) & Treasury yield (same maturity) \\
Interest Coverage Ratio & Senior / subordinated & VIX \\
Net Debt / Equity & Secured / unsecured & IG OAS (BAMLC0A0CM) \\
FCF / Total Debt & Loan size (\$M) & HY OAS (BAMLH0A0HYM2) \\
Revenue Growth (YoY) & Covenant package & GDP growth (YoY) \\
EBITDA Margin & Industry (NAICS) & Yield curve slope (10Y-2Y) \\
Current Ratio & Borrower credit rating & Unemployment rate \\
\bottomrule
\end{tabular}
\end{table}

\subsection{Dataset Summary}

\phtext{Insert dataset summary statistics table here. Include: number of observations, date range, mean/median credit spread, default rate, distribution of loan maturities, and industry breakdown. Example output from \texttt{notebooks/01\_eda.ipynb}.}

\phtext{Insert credit spread distribution histogram here (output of \texttt{notebooks/01\_eda.ipynb}, Figure 1).}

% ─────────────────────────────────────────────────────────────────────────────
\section{Model Training and Evaluation}
\label{sec:training}
% ─────────────────────────────────────────────────────────────────────────────

\subsection{Stage 1: Probability of Default Model}

The first stage trains an XGBoost classifier on borrower financial ratios to predict
a binary default indicator ($y = 1$ if the borrower defaulted within the loan term).

\begin{lstlisting}[caption={Stage 1 -- XGBoost PD classifier (excerpt from \texttt{src/models/pd\_model.py}).}, label={lst:pd}]
# <<< PLACEHOLDER: paste final Stage 1 training code here >>>
import xgboost as xgb
from sklearn.model_selection import StratifiedKFold
from sklearn.metrics import roc_auc_score

pd_model = xgb.XGBClassifier(
    n_estimators=...,       # PLACEHOLDER
    max_depth=...,          # PLACEHOLDER
    learning_rate=...,      # PLACEHOLDER
    scale_pos_weight=...,   # PLACEHOLDER (handles class imbalance)
    eval_metric='auc',
    random_state=42,
)
pd_model.fit(X_train_pd, y_train_pd,
             eval_set=[(X_val_pd, y_val_pd)],
             early_stopping_rounds=50)
\end{lstlisting}

\phtext{Insert Stage 1 evaluation metrics: ROC-AUC (train/val/test), PR-AUC, Brier score, calibration curve. Output of \texttt{notebooks/02\_pd\_model.ipynb}.}

\phtext{Insert SHAP feature importance bar chart for the PD model (top 15 features). Output of \texttt{notebooks/02\_pd\_model.ipynb}, Figure 2.}

\subsection{Stage 2: Credit Spread Regression}

The second stage trains an XGBoost regressor to predict the observed credit spread
$s$ from the estimated PD and additional features.

\begin{lstlisting}[caption={Stage 2 -- XGBoost spread regressor (excerpt from \texttt{src/models/spread\_model.py}).}, label={lst:spread}]
# <<< PLACEHOLDER: paste final Stage 2 training code here >>>
spread_model = xgb.XGBRegressor(
    n_estimators=...,       # PLACEHOLDER
    max_depth=...,          # PLACEHOLDER
    learning_rate=...,      # PLACEHOLDER
    objective='reg:squarederror',
    random_state=42,
)
spread_model.fit(X_train_spread, y_train_spread,
                 eval_set=[(X_val_spread, y_val_spread)],
                 early_stopping_rounds=50)
\end{lstlisting}

\phtext{Insert Stage 2 evaluation metrics: RMSE, MAE, R$^2$, mean absolute percentage error (MAPE) on held-out test set. Scatter plot of predicted vs.\ actual spread. Output of \texttt{notebooks/03\_spread\_model.ipynb}.}

\phtext{Insert SHAP feature importance bar chart and summary beeswarm plot for the spread model. Output of \texttt{notebooks/03\_spread\_model.ipynb}, Figures 3--4.}

\subsection{Hyperparameter Tuning}

\phtext{Describe the hyperparameter search procedure. Include: search method (e.g., Optuna TPE sampler), parameter grid, number of trials, cross-validation strategy, and the final selected hyperparameters for both stages. Summarise in a table. Output of \texttt{notebooks/04\_hyperparameter\_tuning.ipynb}.}

% ─────────────────────────────────────────────────────────────────────────────
\section{Application to Private / Unlisted Borrowers}
\label{sec:private}
% ─────────────────────────────────────────────────────────────────────────────

\subsection{The Challenge}

Public-market models (Merton, KMV) are inapplicable to private borrowers because:
\begin{itemize}
  \item No observable equity price or implied volatility.
  \item No market-traded CDS or bonds from which to extract spreads.
  \item Limited or unaudited financial disclosure.
\end{itemize}

\subsection{Approach 1: Altman Z''-Score}

\citet{altman1995anatomy} adapted the original Z-Score for private firms by replacing
market capitalisation with book equity:

\begin{equation}
  Z'' = 6.56 X_1 + 3.26 X_2 + 6.72 X_3 + 1.05 X_4,
\end{equation}

where $X_1 = \frac{\text{Working Capital}}{\text{Total Assets}}$,
$X_2 = \frac{\text{Retained Earnings}}{\text{Total Assets}}$,
$X_3 = \frac{\text{EBIT}}{\text{Total Assets}}$,
$X_4 = \frac{\text{Book Equity}}{\text{Total Liabilities}}$.

The $Z''$ score is mapped to an internal credit rating and, subsequently, to a
PD and spread.

\subsection{Approach 2: Transfer Learning from Public Firms}

Our two-stage model uses \emph{only accounting-based features}, with no market prices.
It is therefore directly applicable to private borrowers: supply the financial
statements, and the model returns $\widehat{\text{PD}}$ and $\hat{s}$.

\subsection{Approach 3: Internal Ratings-Based (IRB) Scorecard}

Consistent with the Basel III IRB approach, we build a scorecard that assigns the
borrower to an internal rating bucket $r \in \{1, \ldots, 10\}$:

\begin{equation}
  \hat{s}(r) = \overline{s}_r + \beta_1 \cdot \text{Maturity} + \beta_2 \cdot \text{Collateral} + \varepsilon,
\end{equation}

where $\overline{s}_r$ is the historically observed median spread for rating bucket
$r$.
A \textbf{private company illiquidity premium} of approximately 100--200\,bps is
added on top to compensate for the reduced secondary market liquidity of private
loans.

\phtext{Insert worked example: apply the model to a synthetic private borrower with given financial ratios and loan terms. Show the pipeline output: Z''-score, internal rating, estimated PD, estimated spread, and all-in interest rate. Output of \texttt{notebooks/05\_private\_borrower\_example.ipynb}.}

% ─────────────────────────────────────────────────────────────────────────────
\section{Forecasting the Resale Price After One Month}
\label{sec:forecast}
% ─────────────────────────────────────────────────────────────────────────────

\subsection{Loan Mark-to-Market Pricing}

The mark-to-market price of a fixed-rate term loan at time $t$ is:

\begin{equation}
  P(t) = \sum_{k=1}^{n} \frac{C}{(1 + r_f(t) + s(t))^{t_k}} + \frac{F}{(1 + r_f(t) + s(t))^{T}},
\end{equation}

where $C$ is the periodic coupon, $F$ the face value, and $T$ the maturity.
A change in price over one month $[t, t+\Delta t]$ is therefore driven by
(i) changes in the Treasury yield $\Delta r_f$ and
(ii) changes in the credit spread $\Delta s$.

\subsection{Spread Dynamics: Ornstein-Uhlenbeck Model}

We model the credit spread as a mean-reverting Ornstein-Uhlenbeck (OU) process:

\begin{equation}
  ds_t = \kappa(\theta - s_t)\,dt + \sigma\,dW_t,
  \label{eq:ou}
\end{equation}

where $\kappa > 0$ is the mean-reversion speed, $\theta$ is the long-run mean
spread, $\sigma$ is the spread volatility, and $W_t$ is a standard Brownian motion.
The one-month conditional distribution of $s_{t+\Delta t}$ given $s_t$ is Gaussian:

\begin{equation}
  s_{t+\Delta t} \mid s_t \;\sim\; \mathcal{N}\!\left(
    \theta + (s_t - \theta)e^{-\kappa\Delta t},\;
    \frac{\sigma^2}{2\kappa}(1 - e^{-2\kappa\Delta t})
  \right).
\end{equation}

The Treasury yield dynamics are modelled separately via a Vasicek process,
calibrated to the current yield curve.

\subsection{Monte Carlo Forecasting Procedure}

\begin{enumerate}
  \item Estimate OU parameters $(\kappa, \theta, \sigma)$ from historical spread time series via MLE.
  \item Estimate Vasicek parameters for the Treasury yield.
  \item Simulate $N = 10{,}000$ paths of $(r_f(t+\Delta t),\, s(t+\Delta t))$ over
    $\Delta t = 1/12$ year.
  \item For each path, compute the loan price $P(t+\Delta t)$ using Equation~(5).
  \item Report the mean forecast $\hat{P}$ and the empirical 2.5th/97.5th percentiles.
\end{enumerate}

\phtext{Insert Monte Carlo forecast output: histogram of simulated loan prices at $t+1$ month, with mean, 2.5th, and 97.5th percentiles marked. Output of \texttt{notebooks/06\_monte\_carlo\_forecast.ipynb}, Figure 5.}

\phtext{Insert table of OU calibration results: estimated $\hat{\kappa}$, $\hat{\theta}$, $\hat{\sigma}$, along with standard errors and log-likelihood. Output of \texttt{notebooks/06\_monte\_carlo\_forecast.ipynb}.}

% ─────────────────────────────────────────────────────────────────────────────
\section{Uncertainty Quantification: 95\% Confidence Intervals}
\label{sec:ci}
% ─────────────────────────────────────────────────────────────────────────────

\subsection{Static Pricing Interval (Current Fair Spread)}

We implement \textbf{quantile gradient boosting} to produce direct 95\% prediction
intervals for the credit spread without distributional assumptions:

\begin{equation}
  \hat{s}_{\alpha}(x) = \underset{f \in \mathcal{F}}{\arg\min}\;
    \mathbb{E}\!\left[\rho_\alpha\!\left(s - f(x)\right)\right], \quad
    \rho_\alpha(u) = u(\alpha - \mathbf{1}_{u < 0}),
\end{equation}

where $\rho_\alpha$ is the pinball (quantile) loss.
We train three models: $\alpha = 0.025$, $\alpha = 0.5$ (point estimate), and
$\alpha = 0.975$.

\begin{lstlisting}[caption={Quantile regression with XGBoost (excerpt from \texttt{src/models/quantile\_model.py}).}, label={lst:quantile}]
# <<< PLACEHOLDER: paste final quantile model code here >>>
model_low = xgb.XGBRegressor(
    objective='reg:quantileerror',
    quantile_alpha=0.025,
    ...  # PLACEHOLDER: hyperparameters
)
model_high = xgb.XGBRegressor(
    objective='reg:quantileerror',
    quantile_alpha=0.975,
    ...  # PLACEHOLDER: hyperparameters
)
model_low.fit(X_train, y_train)
model_high.fit(X_train, y_train)

ci_lower = model_low.predict(X_new)
ci_upper = model_high.predict(X_new)
\end{lstlisting}

We additionally validate coverage using \textbf{split conformal prediction}
\citep{angelopoulos2023conformal}, which guarantees marginal coverage of exactly
$1 - \alpha = 95\%$ on held-out calibration data.

\phtext{Insert empirical coverage table: for $\alpha \in \{0.90, 0.95, 0.99\}$, report nominal coverage, actual coverage on the test set, and average interval width in basis points. Output of \texttt{notebooks/07\_uncertainty.ipynb}.}

\subsection{One-Month Forecast Interval}

The 95\% forecast interval for the resale price is derived directly from the Monte
Carlo simulation in Section~\ref{sec:forecast}:

\begin{equation}
  \text{CI}_{95\%} = \left[\hat{P}_{0.025},\; \hat{P}_{0.975}\right],
\end{equation}

where $\hat{P}_{\alpha}$ is the $\alpha$-quantile of the empirical distribution of
$P(t + \Delta t)$ across the $N = 10{,}000$ simulated paths.

\phtext{Insert one-month forecast confidence interval: a fan chart showing the median forecast and the 80\%/95\% prediction bands over the month ahead, for a representative sample loan. Output of \texttt{notebooks/06\_monte\_carlo\_forecast.ipynb}, Figure 6.}

% ─────────────────────────────────────────────────────────────────────────────
\section{Conclusion}
\label{sec:conclusion}
% ─────────────────────────────────────────────────────────────────────────────

We have presented an end-to-end quantitative framework for pricing term loan interest
rates. The key contributions are:

\begin{enumerate}
  \item A \textbf{two-stage ML model} (PD classifier + spread regressor) that achieves
    strong out-of-sample accuracy on historical syndicated loan data.
  \item An extension to \textbf{private and unlisted borrowers} via the Altman Z''-Score,
    an IRB scorecard, and transfer learning from public firm financial statements.
  \item A \textbf{one-month resale price forecast} based on a calibrated
    Ornstein-Uhlenbeck spread dynamics model and Monte Carlo simulation.
  \item \textbf{Calibrated 95\% confidence intervals} via quantile gradient boosting,
    validated by split conformal prediction.
\end{enumerate}

\phtext{Insert final model performance summary table: test RMSE, MAE, $R^2$, coverage of 95\% CI, and mean CI width in bps, compared against a benchmark (e.g., OLS with Altman Z-Score). Output of \texttt{notebooks/08\_summary.ipynb}.}

% ─────────────────────────────────────────────────────────────────────────────
\appendix
\section{Mathematical Derivations}
% ─────────────────────────────────────────────────────────────────────────────

\subsection{Merton (1974) Bond Pricing Formula}

Under the Merton framework, the value of a risky zero-coupon bond with face value
$F$ and maturity $T$ is:

\begin{equation}
  D_0 = F e^{-rT} - (F e^{-rT} \Phi(-d_2) - V_0 \Phi(-d_1)),
\end{equation}

where $d_{1,2} = \frac{\ln(V_0/F) + (r \pm \frac{1}{2}\sigma^2)T}{\sigma\sqrt{T}}$
and $\Phi$ is the standard normal CDF.
The credit spread is $s = -\frac{1}{T}\ln(D_0 / Fe^{-rT})$.

\subsection{Ornstein-Uhlenbeck MLE Estimation}

Given discrete observations $\{s_{t_i}\}_{i=0}^{n}$ with spacing $\Delta$, the
conditional log-likelihood is:

\begin{equation}
  \ell(\kappa, \theta, \sigma) = -\frac{n}{2}\ln(2\pi\sigma_\Delta^2)
    - \frac{1}{2\sigma_\Delta^2}\sum_{i=1}^{n}
    \!\left(s_{t_i} - \mu_\Delta(s_{t_{i-1}})\right)^{\!2},
\end{equation}

where $\mu_\Delta(s) = \theta + (s-\theta)e^{-\kappa\Delta}$ and
$\sigma_\Delta^2 = \frac{\sigma^2}{2\kappa}(1 - e^{-2\kappa\Delta})$.

% ─────────────────────────────────────────────────────────────────────────────
\section{Repository Structure}
% ─────────────────────────────────────────────────────────────────────────────

\begin{verbatim}
loan-pricing/
├── data/
│   ├── raw/            # Downloaded from FRED, SEC EDGAR, SBA
│   └── processed/      # Cleaned and feature-engineered datasets
├── notebooks/
│   ├── 01_eda.ipynb
│   ├── 02_pd_model.ipynb
│   ├── 03_spread_model.ipynb
│   ├── 04_hyperparameter_tuning.ipynb
│   ├── 05_private_borrower_example.ipynb
│   ├── 06_monte_carlo_forecast.ipynb
│   ├── 07_uncertainty.ipynb
│   └── 08_summary.ipynb
├── src/
│   ├── data/
│   │   ├── fetch_fred.py
│   │   ├── fetch_sec.py
│   │   └── preprocess.py
│   ├── features/
│   │   └── engineering.py
│   └── models/
│       ├── pd_model.py
│       ├── spread_model.py
│       ├── quantile_model.py
│       └── ou_calibration.py
├── tests/
├── requirements.txt
└── README.md
\end{verbatim}

% ─────────────────────────────────────────────────────────────────────────────
\begin{thebibliography}{99}
\raggedright

\bibitem[Altman(1968)]{altman1968financial}
Altman, E.~I. (1968).
Financial ratios, discriminant analysis and the prediction of corporate bankruptcy.
\textit{Journal of Finance}, 23(4):589--609.

\bibitem[Altman and Saunders(1997)]{altman1995anatomy}
Altman, E.~I. and Saunders, A. (1997).
An analysis and critique of the BIS proposal on capital adequacy and ratings.
\textit{Journal of Banking \& Finance}, 21(11--12):1721--1742.

\bibitem[Angelopoulos and Bates(2023)]{angelopoulos2023conformal}
Angelopoulos, A.~N. and Bates, S. (2023).
Conformal prediction: A gentle introduction.
\textit{Foundations and Trends in Machine Learning}, 16(4):494--591.

\bibitem[Black and Cox(1976)]{blackcox1976}
Black, F. and Cox, J.~C. (1976).
Valuing corporate securities: Some effects of bond indenture provisions.
\textit{Journal of Finance}, 31(2):351--367.

\bibitem[Chen and Guestrin(2016)]{chen2016xgboost}
Chen, T. and Guestrin, C. (2016).
{XGBoost}: A scalable tree boosting system.
In \textit{Proceedings of the 22nd ACM SIGKDD International Conference on Knowledge Discovery and Data Mining}, pages 785--794.

\bibitem[Duffie and Singleton(1999)]{duffie1999modeling}
Duffie, D. and Singleton, K.~J. (1999).
Modeling term structures of defaultable bonds.
\textit{Review of Financial Studies}, 12(4):687--720.

\bibitem[Jarrow et al.(1997)]{jarrowlando1997}
Jarrow, R.~A., Lando, D., and Turnbull, S.~M. (1997).
A Markov model for the term structure of credit risk spreads.
\textit{Review of Financial Studies}, 10(2):481--523.

\bibitem[Jarrow and Turnbull(1995)]{jarrowturnbull1995}
Jarrow, R.~A. and Turnbull, S.~M. (1995).
Pricing derivatives on financial securities subject to credit risk.
\textit{Journal of Finance}, 50(1):53--85.

\bibitem[Lando(1998)]{lando1998cox}
Lando, D. (1998).
On Cox processes and credit risky securities.
\textit{Review of Derivatives Research}, 2(2--3):99--120.

\bibitem[Leland(1994)]{leland1994}
Leland, H.~E. (1994).
Corporate debt value, bond covenants, and optimal capital structure.
\textit{Journal of Finance}, 49(4):1213--1252.

\bibitem[Longstaff and Schwartz(1995)]{longstaffschwartz1995}
Longstaff, F.~A. and Schwartz, E.~S. (1995).
A simple approach to valuing risky fixed and floating rate debt.
\textit{Journal of Finance}, 50(3):789--819.

\bibitem[Mai et al.(2019)]{mai2019deep}
Mai, F., Tian, S., Lee, C., and Ma, L. (2019).
Deep learning models for bankruptcy prediction using textual disclosures.
\textit{European Journal of Operational Research}, 274(2):743--758.

\bibitem[Merton(1974)]{merton1974pricing}
Merton, R.~C. (1974).
On the pricing of corporate debt: The risk structure of interest rates.
\textit{Journal of Finance}, 29(2):449--470.

\end{thebibliography}
% ─────────────────────────────────────────────────────────────────────────────

\end{document}
